%% kapa: exam=1 round=21 year=2010 period=1 form=A subject=economics subject_ko=경제원론 seq=36 number=76
%% source: Q-Net 감정평가사 기출문제 (cst00309) artlSeq=5024181 file=extracted/1008/5024181/2010년 제21회 감정평가사 1차시험 기출문제/1교시.hwp
%% conversion: kordoc 4.24.51 -> tools/gen1.py (auto); review=proofread-sonnet; flags=none
\begin{question}{76}
화폐수량설에 따른 화폐수량방정식은 M × V = P × Y 이다. 이에 관한 설명으로 옳지 않은 것은? (단, M은 통화량, V는 화폐유통속도, P는 산출물의 가격, Y는 산출량이다.)
\begin{choices}
  \item 장기적으로 화폐의 중립성이 성립한다.
  \item 화폐유통속도는 오랜 기간에 걸쳐 비교적 안정적이다.
  \item 중앙은행이 통화량을 증가시키면 물가는 상승한다.
  \item 중앙은행이 통화량을 증가시키면 장기적으로 산출량이 증가한다.
  \item 중앙은행이 통화량을 변화시키면 산출량의 명목가치는 비례적으로 변한다.
\end{choices}
\end{question}
\answer{4}
